% Owner: B. Measured-result-independent methods, PLAN v3.1.
\paragraph{Citation data and time resolution.}
We stream AMiner DBLP-Citation V14 \citep{tang2008arnetminer} and retain an induced subgraph of database
venues selected by a fixed whitelist (SIGMOD, VLDB/PVLDB, ICDE, PODS, EDBT,
ICDT, TODS, VLDB Journal, Data \& Knowledge Engineering, and Information Systems).
All selected papers published in 1960--2020 are retained, including isolates.
A distinct citation $u\to v$ occurs in the publication year of $u$.
The primary analysis requires $y_v<y_u$; same-year citations are saved for
structural sensitivity analysis because their order cannot be resolved.
We exclude self-citations to the same paper, duplicate edges and future references.
Author self-citations are not removed. Counts refer to citations within the
selected venues, rather than global citation impact. The primary analysis years
are 1990--2020, initialized using the available earlier history.

\paragraph{M1: annual and cumulative citation structure.}
Annual and cumulative networks use the same node denominator: all selected
papers already published by year $t$. Directed density is $E/[N(N-1)]$ and
mean total degree is $2E/N$. Clustering, component counts and largest-component
fractions use the weak undirected projection, including isolates.
Trailing windows of 1, 2, 5 and 10 years provide a window-sensitivity check;
sensitivity is $(\max-\min)/\mathrm{mean}$ over the grid at each endpoint,
then averaged over endpoints.

\paragraph{M2c-1: preferential attachment.}
Following the degree-class approach of \citet{jeong2003measuring}, we freeze
in-degree at the start of each citing year and count eligible nodes $n_{tk}$ and
new citations $h_{tk}$ in each prior-degree class $k$.
We estimate $\alpha$ by conditional likelihood with probabilities proportional
to $n_{tk}(k+1)^\alpha$, conditioning on the total number of citations in each year.
The offset admits citations to previously uncited papers; offsets 0.5 and 2
provide sensitivity checks. Confidence intervals use 400 resamples of citing-year
blocks. These are descriptive uncertainty intervals; neither positive $\alpha$
nor its departure from an age-only null identifies a causal mechanism.

\paragraph{M2c-2: ageing.}
For ages 1--60, the target probability in citing year $t$ is proportional to
$n_{ta}f(a)$, where $n_{ta}$ counts available papers of age $a$.
A conditional target likelihood compares discrete-grid lognormal, exponential,
power-law and uniform kernels using AIC on identical observations and support.
We repeat the comparison by publication cohort (1970--1989, 1990--1999,
2000--2009), accounting for each cohort's actual exposure.
This estimates marginal ageing; degree, fitness and venue coverage can also
influence this curve.

\paragraph{M2c-3: densification and diameter.}
We fit $\log E(t)=a\log N(t)+b$ over 1990--2020 and repeat over split periods and
an earlier observation cutoff. Year-resampling intervals are descriptive because
cumulative observations are dependent.
Effective diameter is the 90th percentile of weak shortest-path distances
within the largest connected component, using 128 uniformly sampled source nodes
per year. Sensitivity checks use 64 and 256 sources. We save the component size
and number of sources with every estimate \citep{leskovec2005graphs}.

\paragraph{M2c-4: citation dynamics and prediction.}
We fit the Wang--Song--Barab\'asi curve \citep{wang2013quantifying},
\begin{equation}
 c(t)=m\left[\exp\!\left\{\lambda\Phi\left(\frac{\log t-\mu}{\sigma}\right)\right\}-1\right],
\end{equation}
where $m$ is fixed to the mean observed within-subgraph reference count of papers
published through 1989. Bounded multistart least squares minimizes squared
$\log(1+c)$ errors; parameter bounds and boundary flags are preserved.
For descriptive full-history fits we randomly select at most 500 papers born in
1990--2005 with at least 10 citations by age 15 and fit years 1--15. The resulting
collapse is in-sample evidence only.
For prediction, selection uses only the first five years: at least five observed
citations by age five, followed by a fixed random sample of at most 500 papers.
Using only years 1--5, we compare WSB against no-ageing PA,
$c(t)=m[\exp(rt)-1]$, and linear extrapolation, $c(5)t/5$.
All models predict the same papers at age 15 and at the observed endpoint in 2020;
these are censored observed totals, not unobserved lifetime citations.
We report MAE, median absolute error, RMSLE, paired bootstrap differences
(1000 resamples) and an interior-fit sensitivity check. Boundary fits remain in
the primary evaluation so that difficult cases are not silently discarded.

\paragraph{P3: citation-target randomization.}
Twenty realizations per model preserve each paper's publication year and exact
number of distinct observed references, with no replacement and strictly older
targets. P3a selects eligible targets uniformly. P3b weights targets by a nonparametric age kernel fitted by conditional
target likelihood, accounting for each citing year's available age distribution,
and removes degree preference.
Its age weights are calibrated from the data, so agreement on ageing is not an
independent confirmation. We report null means and sample standard deviations;
empirical tail probabilities have minimum resolution $1/21$.
Both models fix $N(t)$ and $E(t)$ and hence exactly preserve the densification
exponent. They can test diameter and target preference, but cannot supply a
significance test for densification itself.

\paragraph{Reproducibility.}
All significant CPU workloads run on Vanda compute nodes through PBS without GPU
allocation. Fixed seeds, input/source checksums, environment lock file, filtering
counts, individual forecasts and every randomized edge list are saved.
Robustness checks cover time splits, cutoff 2015, degree offsets, same-year edges,
core database venues, diameter sampling, and 10 versus 20 null realizations.
